\chapter{Lezione 7 -- Funzioni pseudocasuali e sicurezza CPA}

\emph{Segue gli appunti del prof. Le parti aggiunte per capire meglio sono segnate ``non negli appunti del prof''. Usa i concetti della Lezione 3 (collision probability, distanza statistica), della Lezione 5 (PRG, ibridi, riduzioni) e della Lezione 6 (cifratura con PRG, game hopping).}

\section{Il problema: cifrare tanti messaggi con la stessa chiave}

\begin{inbreve}
Lo schema della Lezione 6, \(\text{Enc}(k, m) = G(k) \oplus m\), è sicuro solo per \textbf{un} messaggio: con due messaggi e la stessa chiave, \(c_1 \oplus c_2 = m_1 \oplus m_2\). Per cifrare \textbf{tanti} messaggi con \textbf{una sola} chiave corta serve uno strumento nuovo: le \textbf{funzioni pseudocasuali} (PRF).
\end{inbreve}

\section{Funzioni pseudocasuali (PRF)}

\subsection{Definizione informale (del prof)}

\begin{inbreve}
Una \textbf{PRF} è una famiglia di funzioni indicizzata da una chiave \(k\). Scelta la chiave a caso, la funzione \(F_k\) si comporta come una funzione \textbf{completamente casuale}: chi la interroga su degli input e vede gli output non riesce a capire se sta parlando con \(F_k\) o con una funzione scelta a caso tra tutte quelle possibili.
\end{inbreve}

Una PRF è una \textbf{famiglia di funzioni}
\[
\mathcal{F} = \{ F_k : \{0,1\}^n \to \{0,1\}^n \}_{k \in \{0,1\}^\lambda}
\]

\begin{nota}
\textbf{Cosa sono le cose}
\begin{itemize}
  \item \(k\) = la chiave, \(\lambda\) bit. Ogni chiave sceglie una funzione della famiglia: \(F_k\).
  \item \(F_k\) = una funzione che prende \(n\) bit e restituisce \(n\) bit. È \textbf{deterministica}: stesso input, stesso output.
  \item \(R : \{0,1\}^n \to \{0,1\}^n\) = una \textbf{funzione uniformemente casuale}: anche lei deterministica, ma \textbf{scelta a caso tra tutte} le funzioni possibili da \(\{0,1\}^n\) a \(\{0,1\}^n\).
\end{itemize}
\end{nota}

\textbf{Sicurezza (del prof)}: per \(k \leftarrow U_\lambda\) scelta a caso, \textbf{nessun avversario efficiente} \(A\) riesce a distinguere campioni di \(F_k(\cdot)\) da campioni di una funzione uniformemente casuale \(R\).

\begin{esempio}
\textbf{Cos'è una funzione casuale} \emph{(non negli appunti del prof)}. Con \(n = 2\), una funzione da \(\{0,1\}^2\) a \(\{0,1\}^2\) è una \textbf{tabella} con 4 righe:

\begin{center}
\begin{tabular}{cc}
\toprule
\(x\) & \(R(x)\) \\
\midrule
00 & 11 \\
01 & 01 \\
10 & 11 \\
11 & 00 \\
\bottomrule
\end{tabular}
\end{center}

\(R\) si sceglie riempiendo ogni riga con un valore a caso. È deterministica (\(R(10)\) vale sempre 11), ma la tabella è casuale. Le tabelle possibili sono \(4^4 = 256\); in generale sono \((2^n)^{2^n}\), un numero enorme. Una famiglia PRF ne ha solo \(2^\lambda\) (una per chiave): deve ``sembrare'' una tabella casuale pur essendo una delle poche della famiglia.
\end{esempio}

\begin{nota}
\textbf{PRG contro PRF} \emph{(non negli appunti del prof)}: un PRG dà \textbf{una} stringa lunga che sembra casuale. Una PRF dà una tabella enorme (\(2^n\) righe) che sembra casuale, ma non serve calcolarla tutta: si calcola solo la riga che serve, \(F_k(x)\), quando serve. È come un PRG ad ``accesso casuale''.
\end{nota}

\subsection{Definizione formale (del prof)}

\(\mathcal{F}\) è una \textbf{famiglia PRF} se
\[
\text{GAME}^{\text{real}}_{\mathcal{F}, A}(\lambda) \ \approx_c\ \text{GAME}^{\text{ideal}}_{\mathcal{F}, A}(\lambda)
\]

\begin{verbatim}
               GAME^real  /  GAME^ideal
A                                     C
                                      real:  k ← U_λ
                                      ideal: R ← R(n,n)  (tutte le funzioni)
  ┌──────────── x ∈ {0,1}^n ─────────→
  │                                   real:  y = F_k(x)
  │                                   ideal: y = R(x)
  └←─────────── y ∈ {0,1}^n ──────────   (ripetuto quante volte A vuole)
  ──────────────── b' ───────────────→
\end{verbatim}

\begin{itemize}
  \item Nel gioco \textbf{reale} il challenger sceglie una chiave \(k\) a caso e risponde alle domande di \(A\) con \(F_k(x)\).
  \item Nel gioco \textbf{ideale} il challenger sceglie una funzione \(R\) a caso tra \textbf{tutte} quelle da \(\{0,1\}^n\) a \(\{0,1\}^n\) (l'insieme \(\mathcal{R}(n, n)\)) e risponde con \(R(x)\).
  \item \(A\) fa quante domande vuole (\textbf{repeat}), sempre un numero polinomiale perché è efficiente, poi risponde \(b'\).
\end{itemize}

\textbf{Osservazione (del prof)}: il gioco ideale richiede \textbf{tempo esponenziale}, cioè il challenger è illimitato: la tabella di \(R\) ha \(2^n\) righe.

\begin{nota}
\emph{(Non negli appunti del prof.)} Non è un problema, perché il challenger è un oggetto matematico, non un algoritmo che qualcuno deve eseguire. Si può comunque simulare in modo efficiente con il \textbf{lazy sampling}: si sceglie il valore \(R(x)\) a caso solo quando \(A\) chiede \(x\) per la prima volta, e lo si ricorda per le domande successive. \(A\) fa poche domande, quindi la tabella resta piccola.
\end{nota}

\subsection{Esercizio: nessuna PRF è sicura contro avversari illimitati}

\begin{approfondimento}{Soluzione proposta (non negli appunti del prof)}
Come per i PRG (Lezione 5). \(A\) chiede \(F(x_1), \dots, F(x_t)\) su \(t\) input diversi e poi prova \textbf{tutte} le \(2^\lambda\) chiavi: risponde ``reale'' se esiste una chiave \(k\) con \(F_k(x_i) = y_i\) per ogni \(i\).
\begin{itemize}
  \item Nel gioco reale una chiave così esiste sempre (quella vera): \(A\) risponde ``reale'' con probabilità 1.
  \item Nel gioco ideale le risposte sono \(t\) stringhe casuali di \(n\) bit. Una chiave fissata le indovina tutte con probabilità \(2^{-tn}\); per lo union bound, qualche chiave ci riesce con probabilità al più \(2^\lambda \cdot 2^{-tn}\). Con \(t n \ge \lambda + 1\) è \(\le 1/2\).
\end{itemize}
Vantaggio \(\ge 1/2\). Ma provare \(2^\lambda\) chiavi richiede tempo esponenziale: contro un avversario efficiente non funziona.
\end{approfondimento}

\subsection{Come si costruiscono le PRF (del prof)}

\begin{itemize}
  \item \textbf{In pratica}: AES, DES, qualunque \textbf{block cipher}. Nessuno sa \textbf{dimostrare} che AES è una PRF: lo si assume, perché resiste agli attacchi da decenni.
  \item \textbf{In teoria}:
  \begin{itemize}
    \item dalle \textbf{OWF} (si può, ma non è pratico);
    \item dalle stesse ipotesi della crittografia a chiave pubblica (\textbf{fattorizzazione}, \textbf{logaritmo discreto}, \textbf{reticoli}): pratiche ed efficienti, ma comunque meno efficienti di AES.
  \end{itemize}
\end{itemize}

\begin{nota}
\textbf{Block cipher} \emph{(definizione non negli appunti del prof)}: un algoritmo che cifra blocchi di lunghezza fissa (AES: 128 bit) con una chiave. Visto come funzione \(F_k(x)\), cioè ``cifra il blocco \(x\) con la chiave \(k\)'', è il candidato pratico per una PRF.
\end{nota}

\section{Sicurezza per tanti messaggi: CPA}

\begin{inbreve}
Vogliamo uno SKE pratico: \textbf{una} chiave, \textbf{corta}, per \textbf{tanti} messaggi. La definizione di sicurezza giusta è la \textbf{sicurezza CPA} (\emph{Chosen Plaintext Attack}): è il gioco della Lezione 6, ma in più l'avversario può farsi cifrare \textbf{tutti i messaggi che vuole}, prima e dopo la sfida.
\end{inbreve}

\subsection{Definizione (del prof)}

\(\Pi = (\text{Enc}, \text{Dec})\) è \textbf{CPA-sicuro} se
\[
\text{GAME}^{\text{cpa}}_{\Pi, A}(\lambda, 0) \ \approx_c\ \text{GAME}^{\text{cpa}}_{\Pi, A}(\lambda, 1)
\]

dove \(\text{GAME}^{\text{cpa}}_{\Pi, A}(\lambda, b)\), con \(b \in \{0, 1\}\), è:

\begin{verbatim}
A                                       C
                                        k ← K
  ┌───────────────── m ──────────────→
  └←────────── c ← Enc_k(m) ───────────   (domande di cifratura, quante vuole)
  ────────── m0*, m1* ∈ M ───────────→
  ←───────────────── c* ──────────────   c* ← Enc_k(m_b*)   (la sfida)
  ┌───────────────── m ──────────────→
  └←────────── c ← Enc_k(m) ───────────   (altre domande di cifratura)
  ─────────────────── b' ────────────→
\end{verbatim}

\begin{enumerate}
  \item Il challenger sceglie una chiave \(k\).
  \item \textbf{Domande di cifratura}: \(A\) manda messaggi \(m\) a sua scelta e riceve \(c \leftarrow \text{Enc}_k(m)\), quante volte vuole.
  \item \textbf{Sfida}: \(A\) manda \(m_0^*, m_1^*\) e riceve \(c^* \leftarrow \text{Enc}_k(m_b^*)\).
  \item \(A\) può fare altre domande di cifratura.
  \item \(A\) risponde \(b'\).
\end{enumerate}

\begin{nota}
\textbf{Perché ``chosen plaintext''} \emph{(non negli appunti del prof)}: l'avversario sceglie dei messaggi in chiaro (\emph{plaintext}) e ne vede i cifrati. Modella situazioni reali in cui Eve riesce a far cifrare ad Alice messaggi che conosce (es. le intestazioni standard di un protocollo). L'unica differenza con la Lezione 6 sono le domande di cifratura: senza di esse è la sicurezza one-time.
\end{nota}

\subsection{La cifratura deve essere probabilistica (del prof)}

\textbf{Osservazione}: \(\text{Enc}_k(\cdot)\) \textbf{deve essere randomizzata}:
\[
c \leftarrow \text{Enc}_k(m) \quad \equiv \quad c = \text{Enc}_k(m; r), \qquad r \leftarrow \mathcal{R} \ (\text{randomness space})
\]
Di solito \(\text{Dec}_k(\cdot)\) resta deterministica.

\begin{nota}
\textbf{Cosa vuol dire e perché} \emph{(non negli appunti del prof)}: ogni volta che si cifra si usa una moneta nuova \(r\), quindi lo stesso messaggio cifrato due volte dà \textbf{cifrati diversi}. Se Enc fosse deterministica, \(A\) vincerebbe sempre: chiede la cifratura di \(m_0^*\) (domanda di cifratura), riceve \(c_0\), poi nella sfida confronta: se \(c^* = c_0\) è stato cifrato \(m_0^*\), altrimenti \(m_1^*\).
\end{nota}

\section{Uno SKE CPA-sicuro da una PRF}

\begin{inbreve}
Si sceglie ogni volta un punto casuale \(r\), si calcola \(F_k(r)\) e lo si usa come chiave dell'OTP; \(r\) si manda in chiaro insieme al cifrato. Ogni messaggio usa una ``chiave di OTP'' diversa, \(F_k(r)\), e tutte sembrano casuali perché \(F_k\) è una PRF.
\end{inbreve}

\subsection{La costruzione (del prof)}

Con una PRF \(\mathcal{F}\) qualunque:
\[
\text{Enc}(k, m) = (c_1, c_2), \qquad c_1 = r \leftarrow U_n, \qquad c_2 = F_k(r) \oplus m
\]
\[
\text{Dec}(k, (c_1, c_2)) = F_k(c_1) \oplus c_2 = m
\]

\begin{esempio}
\emph{(Valori inventati, non negli appunti del prof.)} \(n = 4\), chiave \(k\) segreta, e supponiamo \(F_k(0110) = 1011\).

\textbf{Alice cifra} \(m = 1100\): sceglie a caso \(r = 0110\), calcola \(F_k(0110) = 1011\) e \(c_2 = 1011 \oplus 1100 = 0111\). Il cifrato è \((0110,\ 0111)\).

\textbf{Bob decifra}: legge \(c_1 = 0110\), calcola \(F_k(0110) = 1011\) (conosce \(k\)) e fa \(1011 \oplus 0111 = 1100 = m\).

Cifrando di nuovo lo stesso \(m\), Alice sceglie un altro \(r\) (es. 1001) e ottiene un cifrato diverso.
\end{esempio}

\begin{nota}
\emph{(Non negli appunti del prof.)} Eve vede \(r\) ma non conosce \(k\), quindi non sa calcolare \(F_k(r)\): per lei è un valore che sembra casuale, e \(c_2\) è come un OTP. Rispetto alla Lezione 6 la chiave \(k\) resta la stessa per tutti i messaggi: cambia \(r\).
\end{nota}

\subsection{Teorema (del prof)}

Se \(\mathcal{F}\) è una famiglia PRF, lo SKE \(\Pi = (\text{Enc}, \text{Dec})\) qui sopra è \textbf{CPA-sicuro}.

\subsection{Dimostrazione: i giochi}

Sia \(G(\lambda, b) \equiv \text{GAME}^{\text{cpa}}_{\Pi, A}(\lambda, b)\). Bisogna mostrare \(G(\lambda, 0) \approx_c G(\lambda, 1)\).

Si considera \(H(\lambda, b)\): \textbf{uguale a} \(G(\lambda, b)\), \textbf{ma} al posto di \(F_k\) si usa una funzione davvero casuale \(R \leftarrow \mathcal{R}(n, n)\).

\begin{center}
\begin{tabular}{lll}
\toprule
 & \(G(\lambda, b)\) (schema vero) & \(H(\lambda, b)\) (con \(R\)) \\
\midrule
Domanda \(m\) & \(c = (r,\ F_k(r) \oplus m)\) & \(c = (r,\ R(r) \oplus m)\) \\
Sfida & \(c^* = (r^*,\ F_k(r^*) \oplus m_b^*)\) & \(c^* = (r^*,\ R(r^*) \oplus m_b^*)\) \\
\bottomrule
\end{tabular}
\end{center}

(con \(r, r^* \leftarrow U_n\) nuovi ogni volta). Il piano \emph{(catena non negli appunti del prof, è la stessa idea della Lezione 6)}:
\[
G(\lambda, 0) \ \approx_c\ H(\lambda, 0) \ \approx\ H(\lambda, 1) \ \approx_c\ G(\lambda, 1)
\]

\subsection{Lemma 1: \texorpdfstring{\(G(\lambda, b) \approx_c H(\lambda, b)\)}{G ≈c H}}

\textbf{Lemma (del prof)}: \(G(\lambda, b) \approx_c H(\lambda, b)\), per ogni \(b \in \{0, 1\}\).

\textbf{Dimostrazione}: per riduzione alla sicurezza della PRF. Se \(A\) distingue \(G\) da \(H\), costruiamo \(B\) che rompe la PRF.

\begin{verbatim}
A                          B                            C (della PRF)
                                                        z = F_k(r)  oppure  R(r)
  ←────── "START" ────────
  ┌────────── m ─────────→
  │                        ┌────── r ← U_n ─────────→
  │                        └←──────── z ─────────────
  └←── c = (r, z ⊕ m) ────
  ─────── m0*, m1* ──────→
                           ──────── r* ────────────→
                           ←──────── z* ─────────────
  ←── c* = (r*, z* ⊕ mb*) ─
  ────────── b' ─────────→
                           ────────── b' ──────────→
\end{verbatim}

\begin{enumerate}
  \item \(B\) fa partire \(A\).
  \item Per ogni domanda di cifratura \(m\): \(B\) sceglie \(r \leftarrow U_n\), lo chiede al suo challenger e riceve \(z\) (che è \(F_k(r)\) oppure \(R(r)\)), poi risponde ad \(A\) con \(c = (r,\ z \oplus m)\).
  \item Per la sfida \(m_0^*, m_1^*\): \(B\) sceglie \(r^*\), riceve \(z^*\) e risponde con \(c^* = (r^*,\ z^* \oplus m_b^*)\).
  \item \(B\) inoltra la risposta \(b'\) di \(A\).
\end{enumerate}

\textbf{Riduzione perfetta}: se il challenger della PRF usa \(F_k\), \(A\) gioca esattamente \(G(\lambda, b)\); se usa \(R\), esattamente \(H(\lambda, b)\). Quindi il vantaggio di \(B\) contro la PRF è uguale a quello di \(A\) tra \(G\) e \(H\). \(\blacksquare\)

\subsection{Lemma 2: \texorpdfstring{\(H(\lambda, 0) \approx H(\lambda, 1)\)}{H0 ≈ H1}}

\begin{inbreve}
Con una funzione davvero casuale \(R\), la sfida è \((r^*,\ R(r^*) \oplus m_b^*)\): se \(R(r^*)\) non è mai stato visto, è un OTP perfetto. L'unico rischio è che una domanda di cifratura abbia usato \textbf{lo stesso} \(r^*\), rivelando \(R(r^*)\). Ma con \(r\) di \(n\) bit scelti a caso, una collisione succede con probabilità trascurabile.
\end{inbreve}

\textbf{Lemma (del prof)}: \(H(\lambda, 0) \approx_s H(\lambda, 1)\) (statisticamente vicini), purché l'attaccante, anche \textbf{illimitato}, faccia al più \(q(\lambda) = \text{poly}(\lambda)\) domande. In particolare \(H(\lambda, 0) \approx_c H(\lambda, 1)\).

\textbf{Intuizione, senza domande di cifratura.} Se \(A\) non fa domande, la sfida è
\[
c^* = (r^*,\ R(r^*) \oplus m_b^*) \qquad r^* \leftarrow U_n
\]
Siccome \(R\) è una tabella casuale, \(R(r^*)\) è una stringa uniforme \(z^*\): la sfida è distribuita come
\[
c^* = (r^*,\ z^* \oplus m_b^*) \qquad r^*, z^* \leftarrow U_n
\]
cioè un \textbf{OTP} (più \(r^*\), che non c'entra con \(b\)). Quindi \(b\) è perfettamente nascosto.

\textbf{Il problema delle domande di cifratura.} Ma \(A\) può fare domande. Cosa succede se \(A\) chiede \(m = 0^n\) e il challenger sceglie per caso \textbf{lo stesso} \(r^*\) della sfida? Allora
\[
c = (r^*,\ R(r^*) \oplus 0^n) = (r^*,\ R(r^*))
\]
e \(A\) vede \(R(r^*)\) in chiaro: la ``chiave'' dell'OTP della sfida è stata rivelata. Ma questo evento capita solo con probabilità trascurabile.

\textbf{In generale.} Sia \textbf{BAD} l'evento ``\(r^*\) coincide con uno degli \(r\) usati nelle domande di cifratura''.
\begin{itemize}
  \item \textbf{Se BAD non succede}: \(H(\lambda, 0) \equiv H(\lambda, 1)\), perché \(R(r^*)\) non compare da nessun'altra parte e si può sostituire con una stringa uniforme \(z^*\) indipendente da tutto: è l'OTP.
  \item \textbf{BAD succede raramente}:
  \[
  \Pr[\text{BAD}] \ \le\ \Pr_{r_i, r_j}[\exists\, i \ne j : r_i = r_j] \ \le\ \sum_{\substack{i, j = 1 \\ i \ne j}}^{q} \Pr[r_i = r_j] \ =\ \sum_{i, j} \text{Col}(U_n) \ \le\ \binom{q}{2} \cdot 2^{-n}
  \]
  Con \(n = n(\lambda) = \text{poly}(\lambda)\): \(\binom{q}{2} \cdot 2^{-n} = \text{poly}(\lambda) \cdot \text{negl}(\lambda) = \text{negl}(\lambda)\).
\end{itemize}

Si può dimostrare che i due fatti insieme implicano
\[
\text{SD}\big(H(\lambda, 0);\ H(\lambda, 1)\big) \ \le\ \Pr[\text{BAD}] \ =\ \text{negl}(\lambda) \qquad \blacksquare
\]

\begin{nota}
\textbf{I passaggi del conto} \emph{(non negli appunti del prof)}:
\begin{itemize}
  \item \(r_1, \dots, r_q\) sono tutti gli \(r\) scelti nel gioco (domande e sfida). BAD implica che due di loro coincidono, quindi \(\Pr[\text{BAD}] \le \Pr[\exists\, i \ne j: r_i = r_j]\).
  \item \textbf{Union bound} (ripasso di probabilità): la probabilità che succeda almeno una delle coincidenze è al più la somma delle probabilità.
  \item \(\Pr[r_i = r_j] = \text{Col}(U_n) = 2^{-n}\): la collision probability dell'uniforme su \(n\) bit (Lezione 3).
  \item Le coppie \(\{i, j\}\) sono \(\binom{q}{2} = \frac{q(q-1)}{2}\), un polinomio se \(q\) lo è.
\end{itemize}
\textbf{Esempio}: con \(q = 1000\) domande e \(n = 128\), \(\Pr[\text{BAD}] \le 500\,000 \cdot 2^{-128} \approx 10^{-33}\).
\end{nota}

\subsection{Conclusione}

\emph{(Non negli appunti del prof.)} Mettendo insieme i due lemmi:
\[
G(\lambda, 0) \ \approx_c\ H(\lambda, 0) \ \approx\ H(\lambda, 1) \ \approx_c\ G(\lambda, 1)
\]
la somma dei vantaggi è \(\text{negl} + \text{negl} + \text{negl} = \text{negl}\): lo schema è CPA-sicuro.

\section{Riepilogo}

\begin{itemize}
  \item Lo schema \(G(k) \oplus m\) è solo one-time; per tanti messaggi con una chiave servono le \textbf{PRF}.
  \item \textbf{PRF}: famiglia \(\{F_k\}\); con \(k\) casuale, \(F_k\) è indistinguibile da una funzione davvero casuale \(R\) per chi la interroga (\(\text{GAME}^{\text{real}} \approx_c \text{GAME}^{\text{ideal}}\)). Nessuna PRF è sicura contro avversari illimitati.
  \item In pratica: AES e i block cipher (non dimostrati). In teoria: dalle OWF, o da fattorizzazione, logaritmo discreto, reticoli.
  \item \textbf{Sicurezza CPA}: il gioco della Lezione 6 più domande di cifratura a scelta, prima e dopo la sfida. Richiede che Enc sia \textbf{randomizzata}.
  \item \textbf{Costruzione}: \(\text{Enc}(k, m) = (r,\ F_k(r) \oplus m)\) con \(r\) casuale; \(\text{Dec}(k, (c_1, c_2)) = F_k(c_1) \oplus c_2\).
  \item \textbf{Teorema}: PRF \(\Rightarrow\) CPA-sicuro. Prova: \(G \approx_c H\) (riduzione perfetta alla PRF) e \(H(\lambda, 0) \approx H(\lambda, 1)\) (OTP, a meno dell'evento BAD di collisione degli \(r\), con \(\Pr[\text{BAD}] \le \binom{q}{2} 2^{-n}\)).
\end{itemize}
